\documentclass[11pt,a4paper]{article}
\usepackage[T1]{fontenc}
\usepackage{amsmath,amssymb,amsthm,mathtools}
\usepackage{geometry}
\usepackage[hidelinks]{hyperref}
\geometry{margin=1in}

\newtheorem{theorem}{Theorem}[section]
\newtheorem{lemma}[theorem]{Lemma}
\newtheorem{proposition}[theorem]{Proposition}
\newtheorem{corollary}[theorem]{Corollary}
\theoremstyle{definition}
\newtheorem{definition}[theorem]{Definition}
\newtheorem{example}[theorem]{Example}

\author{Part II}
\date{Recorded from the lecturer in 2026-2027.}

\title{Principles of Quantum Mechanics}
\begin{document}
\maketitle
\thispagestyle{empty}
\section*{0. Introduction}

\noindent
$\ast$ Recall features from IB QM (Supplementary Notes 1)

\begin{itemize}
    \item Wave-particle duality:
    
    Waves behaving like particles, e.g.\ light quanta,
    \& vice versa (particles behaving like waves).
    
    \item Wavefunction $\psi(x)$ for a particle --- prob.\ density
    $|\psi(x)|^2$.
    
    (Prob.\ intrinsic to theory.)
    
    \item Observables became operators on wavefns:
    \[
        \hat{x} \to x,
        \qquad
        \hat{p} \to -i\hbar\frac{\partial}{\partial x}.
    \]
    \& lack of commutativity, e.g.
    \[
        [\hat{x},\hat{p}] = i\hbar
    \]
    leads to Uncertainty Relation.
    
    \item Schr\"odinger eqn with Hamiltonian $H$ specifies dynamics
    \& eigenvalues of $H$ are allowed energy levels of system.
    
    \item All of this explains H atom structure \& resulting line spectra.
\end{itemize}

\noindent
$\ast$ Aims of this course:

\begin{itemize}
    \item Reformulate QM in a more powerful, flexible, \& useful form
    --- the Dirac formalism (go beyond wavefns!).
\end{itemize}

Previous description appears as particular realisation
\& also leads to alternatives, e.g.
\[
    \tilde{\psi}(p)
    \qquad \text{with} \qquad
    \hat{x} \to +i\hbar\frac{\partial}{\partial p}
    \quad \& \quad
    \hat{p} \to p.
\]

This more general approach allows
\begin{itemize}
    \item Simpler analysis of known problems, e.g.\ oscillator
\end{itemize}

but also lots of new things
\begin{itemize}
    \item \underline{Spin} of particles
    
    \item \underline{Symmetries} (translations, rotations, for example)
    \& conservation laws.
    
    \item \underline{Identical particles} (bosons \& fermions)
    
    \item Framework for quantising more general systems,
    e.g.\ EM field
    
    \item \underline{Perturbation theory}: practical tool for computations ---
    we will not look at too many detailed applications, but we
    will try to keep track of what all this is for --- physical predictions
\end{itemize}

\section*{1. Dirac Formalism}

\subsection*{1.1 Concepts \& Definitions}

A quantum system is described at each instant by a \underline{state}
$|\psi\rangle$ which belongs to a complex vector space $V$;
any non-zero element of this space $V$ is a possible state.

Thus
\[
    |\psi\rangle \ \&\ |\phi\rangle \in V
    \quad\Rightarrow\quad
    \alpha|\psi\rangle+\beta|\phi\rangle\in V
    \qquad
    \text{for any }\alpha,\beta\in\mathbb{C}.
\]

Physically this is a superposition principle leading to wave-like behaviour,
but these states are \underline{not} \underline{wavefns}.
There are also \underline{dual} or conjugate states
$\langle\phi|$ which belong to the dual space $V^\dagger$.

By definition, states \& duals can be combined or paired to give a
complex number
\[
    \underset{\text{``bra''}}{\langle\phi|},
    \quad
    \underset{\text{``ket''}}{|\psi\rangle}
    \ \longmapsto\
    \underset{\text{``bra(c)ket''}}{\langle\phi|\psi\rangle},
    \qquad
    V^\dagger\times V\to\mathbb{C},
\]
with
\begin{align*}
    \langle\phi|
    \bigl(\alpha_1|\psi_1\rangle+\alpha_2|\psi_2\rangle\bigr)
    &=
    \alpha_1\langle\phi|\psi_1\rangle
    +\alpha_2\langle\phi|\psi_2\rangle,\\
    \bigl(\beta_1\langle\phi_1|+\beta_2\langle\phi_2|\bigr)
    |\psi\rangle
    &=
    \beta_1\langle\phi_1|\psi\rangle
    +\beta_2\langle\phi_2|\psi\rangle.
\end{align*}

The spaces $V$ \& $V^\dagger$ come with an
\underline{inner product}, which can be described as a one-to-one
correspondence between the spaces
\[
    V\longleftrightarrow V^\dagger,
\]
with
\[
    |\psi\rangle
    \longleftrightarrow
    \langle\psi|
    =
    \bigl(|\psi\rangle\bigr)^\dagger
    \qquad
    \text{(use same symbol for corresponding states).}
\]

The inner product is
\[
    |\phi\rangle,\ |\psi\rangle
    \longmapsto
    \langle\phi|\psi\rangle
    =
    \bigl(|\phi\rangle\bigr)^\dagger|\psi\rangle,
    \qquad
    V\times V\to\mathbb{C},
\]
\& it is assumed to obey
\[
    \langle\phi|\psi\rangle
    =
    \langle\psi|\phi\rangle^*
    \qquad
    \underline{\text{hermitian}},
\]
\[
    \bigl\|\,|\psi\rangle\,\bigr\|^2
    =
    \langle\psi|\psi\rangle
    \qquad
    \text{real},
\]
\[
    \left.
    \begin{aligned}
        \langle\psi|\psi\rangle &\geq 0,\\
        \langle\psi|\psi\rangle &=0
        \quad\text{iff}\quad
        |\psi\rangle=0
    \end{aligned}
    \right\}
    \qquad
    \underline{\text{positive definite}},
\]
where now $\bigl\|\,|\psi\rangle\,\bigr\|$ is called the norm.

\medskip

\noindent
Note: Knowing $\langle\phi|\psi\rangle$ for all
$\langle\phi|$ determines $|\psi\rangle$
(\& vice versa).

\medskip

\noindent
[Check:
\[
    \langle\phi|\psi_1\rangle
    =
    \langle\phi|\psi_2\rangle
    \quad\Rightarrow\quad
    \langle\phi|\chi\rangle=0,
    \qquad
    \text{where }
    |\chi\rangle=|\psi_1\rangle-|\psi_2\rangle.
\]
If this holds $\forall\,\langle\phi|$, set
$\langle\phi|=\langle\chi|$ \& get
\[
    \langle\chi|\chi\rangle
    =
    \bigl\|\,|\chi\rangle\,\bigr\|^2
    =
    0
    \quad\Rightarrow\quad
    |\chi\rangle=0,
    \qquad
    \text{as claimed.}
\]
]

\medskip

The physical content of any state is unchanged by
\[
    |\psi\rangle\to\alpha|\psi\rangle
    \qquad
    \text{for any }\alpha\in\mathbb{C},\quad\alpha\neq0.
\]

We will usually normalise to make
$\bigl\|\,|\psi\rangle\,\bigr\|=1$, but can still change
\[
    |\psi\rangle\to e^{i\theta}|\psi\rangle.
\]

Relative phases can be important in expressions like
$\alpha|\psi\rangle+\beta|\phi\rangle$.

The space $V$ is \underline{complete}: we assume that appropriate
sequences \& series \underline{converge}.
A complete inner product space is called a
\underline{Hilbert space}.

\medskip

An \underline{operator} $Q$ is a linear map on states
\[
    |\psi\rangle\longmapsto Q|\psi\rangle,
    \qquad
    V\to V,
\]
with
\[
    Q\bigl(\alpha|\psi\rangle+\beta|\phi\rangle\bigr)
    =
    \alpha Q|\psi\rangle+\beta Q|\phi\rangle.
\]

The same operator can be thought of as acting to the left on dual states,
\[
    \langle\phi|\longmapsto\langle\phi|Q,
    \qquad
    V^\dagger\to V^\dagger,
\]
where the right-hand side is defined by
\[
    \bigl(\langle\phi|Q\bigr)|\psi\rangle
    =
    \langle\phi|\bigl(Q|\psi\rangle\bigr)
    =
    \langle\phi|Q|\psi\rangle.
\]

\begin{center}
    \textbf{End of Lecture 1}
\end{center}
\end{document}
